\documentclass[12pt]{article}
\usepackage[utf8]{inputenc}
\usepackage{amsmath, amssymb}
\usepackage{xcolor}
\usepackage{geometry}
\usepackage{hyperref}
\usepackage{fancyhdr}
\usepackage{enumitem}
\usepackage{minted}
\usepackage{booktabs}
\usepackage{tikz}
\usetikzlibrary{shapes, arrows, positioning}

\geometry{margin=1in}
\hypersetup{colorlinks=true, linkcolor=blue, urlcolor=cyan}

\pagestyle{fancy}
\fancyhf{}
\fancyhead[L]{\textbf{\TOPICTITLE}}
\fancyhead[R]{\thepage}

% ------------------------------- 
% Topic Metadata
% ------------------------------- 
\newcommand{\TOPICTITLE}{Image Processing 1, Correlation, & Convolution}
\title{\TOPICTITLE\\\large Study-Ready Notes}
\author{Compiled by Andrew Photinakis}
\date{\today}

\setlength{\headheight}{15pt}

\begin{document}
\maketitle
\tableofcontents
\newpage

\section*{Image Processing 1, Correlation, & Convolution}

\section{Core Terminology}

\begin{enumerate}
    \item Linear Shift Invariant (LSI) System
          \begin{enumerate}
              \item System where output is a weighted sum of inputs (Linearity) and shifting input shifts output by same amount (Shift Invariance)
              \item
          \end{enumerate}
    \item X
          \begin{enumerate}
              \item Y
          \end{enumerate}
\end{enumerate}


\begin{table}[h!]
    \centering
    \renewcommand{\arraystretch}{1.3}
    \begin{tabular}{|p{2.8cm}|p{4.2cm}|p{3cm}|p{2cm}|p{2.5cm}|}
        \hline
        \textbf{Term} & \textbf{Clear Definition (Paraphrased)}                                                                                                           & \textbf{Mathematical Form} & \textbf{Textbook Section} & \textbf{Related Concepts} \\
        \hline

        Linear Shift Invariant (LSI) System
                      & A system where the output is a weighted sum of inputs (linearity) and shifting the input shifts the output by the same amount (shift invariance).
                      & $h(f_1 + f_2) = h(f_1) + h(f_2)$
                      & Sec 3.2
                      & Convolution, Impulse Response                                                                                                                                                                                                          \\
        \hline

        Convolution
                      & A mathematical operation where a kernel (filter) is flipped and slid across an image to compute a weighted sum of local pixels.
                      & $g = f * h$
                      & Eq 3.15
                      & Correlation, LSI                                                                                                                                                                                                                       \\
        \hline

        Correlation
                      & Similar to convolution, but the kernel is not flipped. It measures similarity between an image patch and a kernel template.
                      & $g = f \otimes h$
                      & Eq 3.13
                      & Template Matching                                                                                                                                                                                                                      \\
        \hline

        Separable Filter
                      & A 2D filter kernel that can be decomposed into two 1D kernels (horizontal and vertical), reducing computation time.
                      & $K = v h^{T}$
                      & Eq 3.20
                      & Optimization, Gaussian Filter                                                                                                                                                                                                          \\
        \hline

        Gaussian Filter
                      & A smoothing filter shaped like a bell curve. It is rotationally symmetric and separable. Used for blurring and removing Gaussian noise.
                      & $G(x,y;\sigma)=\frac{1}{2\pi\sigma^2} e^{-\frac{x^2+y^2}{2\sigma^2}}$
                      & Eq 3.24
                      & Smoothing, Separability                                                                                                                                                                                                                \\
        \hline

        Box Filter
                      & A simple smoothing filter where all weights are equal (averaging filter).
                      & $\frac{1}{K^2}$
                      & Sec 3.2.2
                      & Aliasing, Rect Function                                                                                                                                                                                                                \\
        \hline

        Median Filter
                      & A non-linear filter that replaces a pixel with the median value of its neighborhood. Effective for salt-and-pepper noise removal.
                      & $g(i,j)=\mathrm{median}(\text{Neighborhood})$
                      & Sec 3.3.1
                      & Outlier Rejection, Non-linear                                                                                                                                                                                                          \\
        \hline

        Bilateral Filter
                      & An edge-preserving smoothing filter that weights neighbors by both spatial distance and intensity similarity.
                      & $w_{\text{space}} \cdot w_{\text{range}}$
                      & Eq 3.37
                      & Edge Preserving, Denoising                                                                                                                                                                                                             \\
        \hline

        Padding
                      & Strategy for handling image borders when the filter kernel extends beyond the image boundary.
                      & Zero, Wrap, Clamp, Mirror
                      & Fig 3.13
                      & Boundary Effects                                                                                                                                                                                                                       \\
        \hline

        Steerable Filter
                      & A filter whose response at any angle can be computed as a linear combination of a small set of basis filters.
                      & $G_{\theta} = \cos\theta \, G_{0} + \sin\theta \, G_{90}$
                      & Sec 3.2.3
                      & Gradient, Edge Detection                                                                                                                                                                                                               \\
        \hline
    \end{tabular}
    \caption{Summary of Filtering and System Concepts}
\end{table}


\subsection{Linear Filtering}
\begin{itemize}
    \item Goal: Modify image by replacing each pixel with a linear combination of neighboring pixels
    \item Is foundational tool of image processing
    \item Imagine looking at a single pixel in an image
          \begin{enumerate}
              \item
          \end{enumerate}
\end{itemize}






% template
\begin{itemize}
    \item X
          \begin{enumerate}
              \item Y
          \end{enumerate}
\end{itemize}

\end{document}